% page 386 of the facsimile (image p-385 of the scan)
% block 157.5 x 216.0 mm ; left margin 8.0 mm
\pgc{-12.700mm}{0.000mm}{
\l{\cel{3}378}
\l{}
\l{\sou{modo}. I. Maniero vidar, agar, partikulare di persono, di}
\l{\cel{3}lando, di regiono. - II. Kustumo efemera qua normizas}
\l{\cel{3}(segun la gusto di la periodo) la maniero vivar, vesti-}
\l{\cel{3}\sou{zar}\cel{1}su, e c. - III. Chapeli e kapo-vesti por la homini.}
\l{\cel{3}- IV. (\sou{filoz.}) Eso-maniero variiva di la substanco fi-}
\l{\cel{3}nita. - V. (\sou{logiko)}\cel{2}Kondiciono specala a qua subord-}
\l{\cel{3}inesas silogismo. Dispozeso-maniero di la tri propozi-}
\l{\cel{3}cioni di silogismo segun lia quanteso o qualeso. -VI.}
\l{\cel{3}(yuro-cienco)Klauzo qua subordinas la efiko da ago ad}
\l{\cel{3}evento necerta qua dependas de la volo di ta qua profi-}
\l{\cel{3}tos de olu. - VII. (\sou{muziko)}. 1. Singla de la toni, che}
\l{\cel{3}la muziko antiqua, quin distingas la diapazono plu o}
\l{\cel{3}min alta e la plaso di la mi-toni en la gamo. 2. Singla}
\l{\cel{3}de la toni di la plan-kanto quin distingas (ultre la}
\l{\cel{3}diapazono plu o min alta e la plaso di la mi-toni en la}
\l{\cel{3}gamo) la divideso, en singla gamo, di la oktavo aden}
\l{\cel{3}quarto plasizita altre. - VIII. (\sou{gram.}) Formo di la ver-}
\l{\cel{3}bo, qua indikas la manieri diversa ye qui onu afirmas}
\l{\cel{3}la ago e la stando quin expresas la verbo (exter la cir-}
\l{\cel{3}konstanci di tempo e di personi quin indikas altra formi).}
\l{\cel{3}- DEFIRS.}
\l{}
\l{\sou{modelo}. I.To quo destinesas esor objekto di imito. II. Homo}
\l{\cel{3}qua esas koram artisto (piktisto, statuifisto) en la pos-}
\l{\cel{3}turo quan ica volas reproduktar. - III. Reprezenturo ek}
\l{\cel{3}argilo, vaxo, e c. di subjekto destinata a reproduktesor}
\l{\cel{3}ek marmoro, bronzo, e c. - IV. To quo destinesas esor}
\l{\cel{3}objekto di imito, segun quo onu facas od agas ulo :}
\l{\cel{3}(1) kompozo-maniero ; (2) vivo-maniero. -\cel{2}DEFIRS.}
\l{}
\l{\sou{moderar.}\cel{2}(\sou{trans.}) Forigar de eceso.\cel{2}- DEFIS.}
\l{}
\l{\sou{moderna}. Qua apartenas a nia epoko od ad epoko kompare recen-}
\l{\cel{3}ta (opoze a la epoki olima od antiqua). - DEFIS.}
\l{}
\l{\sou{modesta}. I. For luxego, brilo, pompo. - II. Qua opinionas}
\l{\cel{3}sua merito kom nur mezvalora, pasableta. - EFIS.}
\l{}
\l{\sou{modifikar}. (\sou{trans.}) Chanjar (kozo) sen alterar olua naturo}
\l{\cel{3}esencala. - DEFIRS.}
\l{}
\l{\sou{modiliono}. (\sou{arkitekt.)}\cel{2}Sorto di konsolo, qua havas la formo}
\l{\cel{3}di voluto duopla, salianta, sub la pluv-gutifilo di la kor-}
\l{\cel{3}nico. - DEFIS.}
\l{}
\l{\sou{modlar}. (trans.) Fasonar ek argilo, vaxo, e c. (modelo qua}
\l{\cel{3}exekutesos ye gipso, marmoro, bronzo, e c.). -(\sou{Anke alu-}}
\l{\cel{3}\sou{dante la materio quan onu fasonas)}. - DEFIRS.}
}
